\section{Implementaci\'on y herramientas}

La implementaci\'on se desarroll\'o en Python 3.11 y qued\'o organizada en tres
notebooks finales, ejecutados en orden:

\begin{enumerate}
    \item \texttt{markowitz\_vs\_benchmark\_rebalanceo\_pandemia.ipynb}:
    optimiza carteras Markowitz, calcula el benchmark top-20, aplica rebalanceo
    semestral y compara escenarios con y sin pandemia.

    \item \texttt{black\_litterman.ipynb}: construye portafolios
    Black-Litterman usando las mismas m\'etricas y perfiles, con views de
    momentum, desempleo, momentum top-20 a 6 meses y momentum top-20/bottom-20
    a 1 a\~no.

    \item \texttt{p4\_montecarlo\_recomendador\_portafolios.ipynb}: usa los
    KPIs anteriores como insumo y simula 5\,000 trayectorias para las 180
    combinaciones evaluadas.
\end{enumerate}

El c\'odigo computacional se encuentra disponible en el
\href{https://github.com/Psacevedo/Capstone.git}{repositorio GitHub del proyecto}.
El Anexo N°4 documenta la estructura de estos notebooks, los archivos CSV
generados y la relaci\'on entre c\'odigo, supuestos, KPIs y resultados
experimentales.

Las principales herramientas fueron \textbf{Pandas} para manejo de datos y
exportaci\'on a CSV, \textbf{NumPy/SciPy} para c\'alculo num\'erico y
simulaci\'on, \textbf{CVXPY} para resolver los problemas de optimizaci\'on,
\textbf{Matplotlib} para gr\'aficos, y \textbf{Hashlib} para semillas
deterministas en Monte Carlo.

La trazabilidad se asegura mediante carpetas \texttt{outputs/} por notebook.
Cada etapa exporta CSV con pesos, KPIs, validaciones, rankings o simulaciones,
y esos archivos son los que alimentan las tablas y figuras del informe. Adem\'as,
los notebooks incluyen checklists finales que verifican archivos esperados,
conteos por escenario/t\'ecnica/perfil y ausencia de \texttt{NaN} en columnas
cr\'iticas. 

\section{Pipeline de la Entrega 3}

La tercera fase organiza el pipeline en torno a la calibraci\'on independiente
de cada \textit{view} y a una \'unica regla de abandono. Para cada una de las
cuatro \textit{views} se ejecuta el script
\texttt{validacion\_<view>\_v1\_plus.py} (en la carpeta
\texttt{views\_v1\_plus/}), que internamente:

\begin{enumerate}
    \item \textbf{Calibra} la \textit{view} de forma independiente en cinco
    etapas sobre $\mathcal{H}_1$, con validaci\'on en $\mathcal{H}_2$,
    congelando su propia configuraci\'on de hiperpar\'ametros BL.

    \item \textbf{Valida} con ventanas rolling de tres horizontes, generando
    m\'etricas de Sharpe, drawdown y din\'amica de portafolio (turnover,
    $N_{\text{eff}}$, HHI) por ventana, perfil y escenario.

    \item \textbf{Simula} la evaluaci\'on econ\'omica P4 en el horizonte de test
    limpio $\mathcal{H}_3$, con 2\,000 trayectorias por combinaci\'on, abandono
    semanal (\texttt{views\_v1\_plus}) y utilidad de FinPUC por comisi\'on
    mensual de 0,5\% sobre el saldo activo administrado.
\end{enumerate}

Finalmente, \texttt{generar\_informes\_views\_v1\_plus.py} consolida los
\textit{outputs} de las cuatro \textit{views}, comparando cada metodolog\'ia
Black-Litterman calibrada contra el caso base Markowitz y produciendo las
tablas y figuras agregadas (mejora de Sharpe y de Score P4 ---en promedio y
mediana---, abandono, clientes retenidos y utilidad por saldo).

Todos los scripts utilizan las mismas librer\'ias que la Entrega 2 (Pandas,
NumPy/SciPy, CVXPY, Matplotlib, Hashlib), y los \textit{outputs} se organizan en
carpetas \texttt{outputs/} por m\'odulo, manteniendo la trazabilidad del
proyecto.
